\section*{\S~8a Umgang mit Beschwerden gegen Einzelpersonen}
\label{sec:umgang_mit_beschwerden_gegen_einzelpersonen}
\addcontentsline{toc}{section}{\S~8a Umgang mit Beschwerden gegen Einzelpersonen}

\begin{enumerate}[label=(\arabic*)]
    \item Jedes Mitglied hat, unabhängig von Dienstgrad und Funktion, wenn nicht anders festgelegt, das Recht, Beschwerde gegen Einzelpersonen einzureichen. Gegen jedes Mitglied kann, unabhängig von Dienstgrad und Funktion, Beschwerde eingereicht werden.    
    \item Möchte ein Mitglied in einem Mannschaftsdienstgrad Beschwerde gegen eine Einzelperson einreichen, so ist dies beim Gefreiten vom Dienst (GvD) zu tun. Möchte ein Mitglied in einem Nicht-Mannschaftsdienstgrad Beschwerde gegen eine Einzelperson einreichen, so ist dies beim Personal-Feldwebel zu tun. Ist der Posten des Personal-Feldwebels nicht besetzt, dann ist dies beim Spieß zu tun.    
    \item Wird eine Beschwerde eingereicht, so tagt der Stab nach Anhörung aller Seiten und entscheidet, wie mit der Beschwerde umzugehen ist. Wird diese Entscheidung innerhalb von 72 Stunden nach Bekanntgabe von einem Mitglied angefochten, so tagt der erweiterte Stab nach erneuter Anhörung aller und legt fest, wie mit der Beschwerde umzugehen ist. Die Entscheidung gilt dann endgültig. Gefreiter vom Dienst (GvD) und Personal-Feldwebel vertreten dabei in jedem Fall die Interessen der Mannschaft. Gibt es keinen erweiterten Stab, so tagt der Stab erneut. Eine Entscheidung muss in jedem Fall öffentlich bekannt gemacht werden-    
    \item Geht das Beschwerde einreichende Mitglied davon aus, dass die für die Beschwerde zuständige Stelle diese Beschwerde nicht nachverfolgen kann, so ist die nächst höhere Stelle zuständig. Die Reihenfolge der zuständigen Stellen wird in \S~8a Absatz 4\ref{itm:reihenfolge_beschwerde} festgelegt.
    \begin{enumerate}[label=(\Alph*)]
        \item\label{itm:reihenfolge_beschwerde} Gefreiter vom Dienst (GvD) —> Personal-Feldwebel —> Spieß —> Unit-Leiter —> S1-Feldwebel
        \item Richtet sich die Beschwerde gegen den Unit-Leiter, oder den Spieß, so sind diese beiden Stellen zu überspringen. Vor der öffentlichen Konfrontation gegenüber der entsprechenden Person während einer Nachbesprechung, oder während der Mitgliederversammlung, ist der S1-Feldwebel die höchste anzurufende Stelle, bei Beschwerden solcher Art. Ist der Posten des S1-Feldwebels nicht besetzt, so kann neben dem Gefreiten vom Dienst (GvD) zusätzlich der Unteroffizier vom Dienst (UvD) angerufen werden. Die öffentliche Konfrontation gegenüber der entsprechenden Person während einer Nachbesprechung, oder während der Mitgliederversammlung kann in einem solchen Fall, auch ohne vorheriges Durchlaufen aller Stellen, ergriffen werden. Richtet sich die Beschwerde gegen eine andere Person, die einen der Posten aus \S~8a Absatz 4A inne hat, aber nicht den Posten Unit-Leiter, oder Spieß, so wird diese Stelle übersprungen.
    \end{enumerate}    
    \item Die Stelle, an der Beschwerde eingereicht wurde, hat diese in das kleinste Organ, dem es angehört, einzubringen. Hierbei dürfen nur der Stab und der erweiterte Stab berücksichtigt werden. Dieses wendet auf die Beschwerde dann \hyperref[sec:umgang_mit_regelverstoeßen]{\S~8} an. Ist kein Regelverstoß festzustellen, gilt die Beschwerde als abgelehnt. Wird ein Regelverstoß festgestellt, so greift \S~8 Absatz \ref{itm:vorgehen_bei_verstoß}. Der Person, über die Beschwerde eingereicht wurde, ist dies unverzüglich nach Feststellung eines Regelverstoßes anonymisiert mitzuteilen. Wird diese Entscheidung binnen 72 Stunden angefochten, greift \S~8 Absatz \ref{itm:vorgehen_bei_verstoß} Satz 10.
    \begin{enumerate}[label=(\Alph*)]
        \item Eine einmalige Beschwerde, bei der ein Regelverstoß festgestellt wird, zieht, sofern nicht bereits eine Verwarnung ausgesprochen wurde, eine mündliche Verwarnung nach sich. Diese gilt für 3 Monate.
    \end{enumerate}    
    \item Jedes Mitglied hat das Recht bei der Leitung zu erfahren, wie viele Beschwerden gegen es eingegangen sind, ob aktuell eine Beschwerde vorliegt und bei wie vielen dieser Beschwerden ein Regelverstoß festgestellt wurde.
\end{enumerate}
